% !TeX program = xelatex
% XA-202612「OS Agent 记忆优化及高效应用研究」技术方案报告
% 编译：latexmk -xelatex main.tex（或连续运行两遍 xelatex main.tex 生成目录）
\documentclass[zihao=-4,a4paper]{ctexart}

\usepackage[margin=2.5cm]{geometry}
\usepackage{amsmath}
\usepackage{booktabs}
\usepackage{longtable}
\usepackage{array}
\usepackage{xcolor}
\usepackage{tikz}
\usetikzlibrary{positioning,arrows.meta}
\usepackage{listings}
\usepackage{enumitem}
\usepackage{caption}
\usepackage{url}
\usepackage[hidelinks]{hyperref}

% 段落断行弹性，吸收少量中英混排超宽
\emergencystretch=3em

% ---------- 颜色与 TODO 标记 ----------
% 约定：所有未实测 / 待完成内容统一以红色 TODO 标出（\todo、\todoblank、\statustodo）
\definecolor{todored}{RGB}{211,47,47}
\definecolor{codebg}{RGB}{246,246,246}
\definecolor{codeframe}{RGB}{176,190,197}
\definecolor{figfill}{RGB}{239,246,255}

\newcommand{\todo}[1]{\textcolor{todored}{\textbf{【TODO：#1】}}}
\newcommand{\todoblank}{\textcolor{todored}{\textbf{TODO}}}
\newcommand{\statustodo}{\textcolor{todored}{\textbf{待完成}}}
\newcommand{\statusdone}{已实现}
\newcommand{\statuspartial}{部分实现}

% 行内代码（text 模式转义 \_ \{ 直接正确渲染；不引入断行）
\newcommand{\code}[1]{\texttt{#1}}
% 可断行的长测试名（\path 允许在下划线处换行，参数内下划线无需转义）
\newcommand{\testname}[1]{{\ttfamily\footnotesize\path{#1}}}

\lstdefinestyle{amalst}{
  basicstyle=\ttfamily\footnotesize,
  backgroundcolor=\color{codebg},
  frame=single,
  rulecolor=\color{codeframe},
  breaklines=true,
  columns=fullflexible,
  showstringspaces=false,
  xleftmargin=2mm,
}
\lstset{style=amalst}

\captionsetup{font=small, labelfont=bf}

% TikZ 通用样式
\tikzset{
  box/.style={draw, rounded corners=2pt, align=center, inner sep=5pt,
    font=\small, minimum height=8.5mm, fill=figfill},
  proc/.style={box, fill=gray!10},
  todo/.style={box, dashed, draw=todored, text=todored, fill=white},
  arrow/.style={-{Stealth[length=2.6mm]}, semithick},
}

\begin{document}

% ---------- 封面 ----------
\begin{titlepage}
  \centering
  \vspace*{20mm}
  {\heiti\zihao{2} OS Agent 记忆优化及高效应用研究\par}
  \vspace{8mm}
  {\zihao{3} 多源融合偏好与知识记忆优化技术方案报告\par}
  \vspace{6mm}
  {\zihao{-4} 题目编号：XA-202612\par}
  \vspace{4mm}
  {\zihao{-4} 发榜单位：麒麟软件有限公司\par}
  \vspace{22mm}
  \begin{tabular}{rl}
    参赛单位： & \todo{学校 / 单位全称} \\
    队伍名称： & \todo{队伍名称}         \\
    队伍成员： & \todo{成员姓名}         \\
    指导教师： & \todo{指导教师姓名}     \\
  \end{tabular}
  \vfill
  {\zihao{-4} 2026 年 9 月\par}
  \vspace*{8mm}
\end{titlepage}

\tableofcontents
\newpage

\section{引言与需求分析}\label{sec:intro}

\subsection{题目背景}

基于麒麟操作系统的 OS Agent 是智能交互系统的核心组件，其记忆模块中的\textbf{偏好记忆}
与\textbf{知识记忆}是支撑个性化服务适配、知识沉淀复用的核心载体。当前两类记忆面临的
问题具有多源性。一方面，\textbf{工具调用执行结果结构化不足、数据质量参差不齐}，导致
偏好提取（如工具选择偏好）不精准、知识沉淀不高效；另一方面，\textbf{其他数据源的
问题相互叠加}——用户行为捕捉不全面、跨场景数据不一致、配置版本混乱——最终形成
动态偏好提取偏差、版本化管理效率低、新旧知识冲突滞后、关联检索精度有限等痛点。

本方案针对上述痛点，面向麒麟操作系统 OS Agent 交付一套\textbf{多源融合偏好与知识记忆
优化解决方案}，具备偏好精准捕捉、知识智能整合、高效检索复用等核心功能，并满足端侧
轻量化部署与数据安全合规要求。

\subsection{需求分析}

赛题聚焦两大核心问题，本方案将其分解为五项技术需求。（1）\textbf{偏好记忆的动态捕捉
与适配}：以工具调用执行结果为核心数据源，结合用户行为数据与手动配置，实现用户操作
习惯、输出风格、安全策略等偏好的精准提取、版本化更新与跨场景复用；（2）\textbf{知识
记忆的结构化整合}：新旧知识冲突处理、关联检索优化，支持工作流程知识、历史案例、
可复用模板的结构化存储与智能调用；（3）\textbf{端侧部署约束}：调用银河麒麟国产桌面
操作系统 embedding SDK，轻量化存储，检索响应延迟 $\le 500\mathrm{ms}$；
（4）\textbf{安全合规}：敏感信息识别与过滤，自然语言指令驱动的精准遗忘；
（5）\textbf{记忆流转}：与短期、中期记忆间的数据流转兼容。

\subsection{赛题要求覆盖总览}\label{sec:coverage}

表~\ref{tab:coverage-func} 给出赛题七项内容要求与本方案的对应关系及当前状态；
表~\ref{tab:coverage-metric} 给出四项性能指标的覆盖状态。\textbf{所有未实测 / 待完成
项均以红色 TODO 标出}，详见第~\ref{sec:evaluation}~章效果验证报告。

\begin{table}[htbp]
  \centering\small
  \caption{赛题内容要求覆盖对照}
  \label{tab:coverage-func}
  \begin{tabular}{c p{5.7cm} p{5.2cm} c}
    \toprule
    编号 & 赛题内容要求 & 本方案对应模块 & 状态 \\
    \midrule
    (1) & 多源数据整合（工具结果 / 用户行为 / 手动配置，清洗、标准化、质量校验） &
      三路偏好提取器 + 隐私脱敏边界（\S\ref{sec:pref}） & \statuspartial \\
    (2) & 偏好记忆动态捕捉（自动提取、版本化管理、跨场景适配与回溯） &
      偏好版本链 + 场景标签机制（\S\ref{sec:pref}） & \statuspartial \\
    (3) & 知识记忆结构化整合（冲突处理、关联检索、模板沉淀） &
      SPO 事实库 + 六动作冲突处理 + 模板生命周期（\S\ref{sec:knowledge}） & \statusdone \\
    (4) & 端侧部署（麒麟 embedding SDK、轻量存储、延迟 $\le$ 500ms） &
      可插拔嵌入后端 + int8 量化向量库（\S\ref{sec:embedding}、\S\ref{sec:deploy}） & \statuspartial \\
    (5) & 敏感信息识别过滤 + 自然语言精准遗忘 &
      六类检测器 + 两阶段遗忘（\S\ref{sec:forget}） & \statusdone \\
    (6) & 与短期 / 中期记忆的数据流转 &
      压缩退役网关 + 三层记忆架构（\S\ref{sec:flow}） & \statuspartial \\
    (7) & 量化评测机制与测试报告 &
      四指标评测框架（\S\ref{sec:evaluation}） & \statustodo \\
    \bottomrule
  \end{tabular}
\end{table}

\begin{table}[htbp]
  \centering\small
  \caption{性能指标覆盖对照}
  \label{tab:coverage-metric}
  \begin{tabular}{l c l}
    \toprule
    性能指标 & 目标值 & 本方案现状 \\
    \midrule
    偏好提取准确率 & $\ge 85\%$ & 算法已实现，量化实测 \todoblank \\
    知识检索召回率 & $\ge 85\%$ & 检索机制已实现，量化实测 \todoblank \\
    知识检索响应时间 & $\le 500\mathrm{ms}$ & 线性扫描复杂度分析见 \S\ref{sec:embedding}，实测 \todoblank \\
    知识冲突处理正确率 & $\ge 88\%$ & 六动作机制已实现，量化实测 \todoblank \\
    \bottomrule
  \end{tabular}
\end{table}

\subsection{方案核心思路}

本方案的设计思路可概括为四点。其一是\textbf{确定性优先的记忆算法}：偏好提取与知识
冲突处理均采用可解释的规则与统计阈值（无 LLM 依赖），行为可审计、结果可复现，契合
端侧资源约束。其二是\textbf{只追加的版本化}：偏好与知识均以版本链存储，回滚与回溯是
非破坏性操作，天然支持跨场景适配与按时间点回溯。其三是\textbf{隐私内建}
（privacy by construction）：所有记忆入库边界统一经过敏感信息脱敏，遗忘操作采用
「预览—人工确认—物理删除—残留复验」两阶段协议。其四是\textbf{可插拔嵌入与轻量
存储}：嵌入后端按配置在 remote / local / kylin 三者间切换，向量库采用 int8 量化与
单文件 JSON 持久化，无外部服务依赖。

\subsection{报告组织}

第~\ref{sec:arch}~章总体架构；第~\ref{sec:algo}~章核心算法设计；第~\ref{sec:flow}~章
记忆流转机制；第~\ref{sec:impl}~章技术实现方案；第~\ref{sec:deploy}~章端侧部署与
银河麒麟适配；第~\ref{sec:evaluation}~章效果验证报告；第~\ref{sec:case}~章真实场景
应用案例；第~\ref{sec:conclusion}~章总结与优化方向。

\section{总体架构设计}\label{sec:arch}

\subsection{设计原则}

架构设计遵循四项原则。\textbf{无 LLM 依赖的记忆内核}：偏好提取、冲突处理、遗忘判定
均为确定性规则与统计阈值，不引入额外推理开销，适合端侧部署，且每一步决策可解释、
可审计。\textbf{单一职责的 crate 划分}：领域类型、存储、服务逻辑、接入管道分层清晰，
便于按需裁剪与移植。\textbf{文件即数据库}：全部记忆状态持久化为工作区 \code{.amadeus/}
目录下的 JSON 文件，零外部服务依赖，便于备份、检视与迁移。\textbf{安全边界前置}：
敏感信息检测与脱敏发生在一切入库边界，而非事后清理。

\subsection{三层记忆总体架构}\label{sec:arch-tiers}

系统沿用「短期—中期—长期」三层记忆模型（图~\ref{fig:tierflow}）：

\begin{figure}[htbp]
  \centering
  \begin{tikzpicture}[node distance=7mm and 8mm]
    \node[box] (ctx) {短期记忆\\ {\scriptsize 会话上下文窗口}};
    \node[box, right=of ctx] (mid) {中期记忆库\\ {\scriptsize mid\_term\_memory.json}};
    \node[box, right=of mid] (long) {长期记忆\\ {\scriptsize memory.json}\\[-1pt] {\scriptsize rag\_index.json}};
    \node[proc, right=of long] (prompt) {提示词注入\\ {\scriptsize Persistent Memory}};
    \node[proc, below=10mm of ctx] (multi) {多源输入：工具执行结果 / 用户行为 / 手动配置};

    \draw[arrow] (ctx) -- node[above,font=\scriptsize] {网关} (mid);
    \draw[arrow, dashed, todored] (mid) -- node[above,font=\scriptsize,text=todored] {晋升} (long);
    \draw[arrow] (long) -- node[above,font=\scriptsize] {注入} (prompt);
    \draw[arrow] (prompt.south) to[bend right=26] node[below,font=\scriptsize] {增强上下文} (ctx.south);
    \draw[arrow] (multi) -- (ctx);
  \end{tikzpicture}
  \caption{三层记忆架构与数据流转（短中期之间经 ContextGate 网关；红色虚线为待实现的晋升路径）}
  \label{fig:tierflow}
\end{figure}

三层记忆的分工如下。\textbf{短期记忆}是活跃会话的上下文窗口（\code{Agent.history}）
与压缩摘要，生命周期为单会话；\textbf{中期记忆}是跨会话的记录数据库，沉淀任务、
决策、文件状态、错误与解决方案、状态快照（\S\ref{sec:flow}）；\textbf{长期记忆}则
涵盖用户事实（\code{memory.json}）、结构化知识库（\code{structured\_memory.json}）、
偏好库（\code{preferences.json}）与语义索引（\code{rag\_index.json}），经
\code{MemoryRegistry} 注入系统提示词。

\subsection{双库架构：偏好库与知识库}\label{sec:arch-dual}

偏好与知识两类记忆分别由独立仓库管理，共享统一的隐私脱敏边界（图~\ref{fig:dualrepo}）。

\begin{figure}[htbp]
  \centering
  \begin{tikzpicture}[node distance=6mm and 9mm]
    \node[box] (tool) {工具执行结果};
    \node[box, right=8mm of tool] (behavior) {用户行为数据};
    \node[box, right=8mm of behavior] (config) {手动配置信息};

    \node[proc, below=9mm of behavior] (redact) {统一脱敏与质量校验边界\\
      {\scriptsize SensitiveDataDetector / 输入校验}};

    \node[box, below=10mm of redact, xshift=-24mm] (pref) {偏好库 preferences.json\\
      {\scriptsize 三路提取器 + 版本链}};
    \node[box, below=10mm of redact, xshift=24mm] (know) {知识库 structured\_memory.json\\
      {\scriptsize 事实版本链 + 关联图 + 模板}};

    \node[todo, below=9mm of pref] (pinject) {偏好注入提示词};
    \node[proc, below=9mm of know] (ktool) {kylin\_memory 工具\\ {\scriptsize 查询 / 冲突 / 遗忘}};

    \draw[arrow] (tool) -- (redact);
    \draw[arrow] (behavior) -- (redact);
    \draw[arrow] (config) -- (redact);
    \draw[arrow] (redact) -- (pref);
    \draw[arrow] (redact) -- (know);
    \draw[arrow, dashed, todored] (pref) -- (pinject);
    \draw[arrow] (know) -- (ktool);
  \end{tikzpicture}
  \caption{多源数据接入与偏好 / 知识双库（红色虚线框为待接线部分）}
  \label{fig:dualrepo}
\end{figure}

\subsection{代码组织}

记忆子系统在 Rust 工作区内按职责划分为独立 crate（表~\ref{tab:crates}）。

\begin{table}[htbp]
  \centering\small
  \caption{记忆子系统代码组织}
  \label{tab:crates}
  \begin{tabular}{l l l}
    \toprule
    crate & 职责 & 关键内容 \\
    \midrule
    \code{memory-domain} & 领域类型 & 事实版本、片段、关联边、模板的纯数据模型 \\
    \code{memory-service} & 知识库服务 & 冲突处理、关联检索、模板生命周期、两阶段遗忘 \\
    \code{memory} & 中期记忆 & 记录模型、存储接口、上下文网关（ContextGate） \\
    \code{context}（preference 模块） & 偏好库 & 三路提取器、版本化、场景标签与回溯 \\
    \code{privacy} & 隐私保护 & 六类敏感信息检测器与脱敏 \\
    \code{rag} & 语义检索 & 可插拔嵌入后端、int8 量化向量库、RagTool \\
    \code{core}（memory\_wiring 等） & 接线 & agent 主循环与记忆系统间的数据管道 \\
    \bottomrule
  \end{tabular}
\end{table}

\subsection{运行时存储布局}

全部状态位于工作区 \code{.amadeus/} 目录（表~\ref{tab:storage}），单文件 JSON、
人可直接检视，符合端侧轻量化与可审计要求。

\begin{table}[htbp]
  \centering\small
  \caption{运行时存储文件}
  \label{tab:storage}
  \begin{tabular}{l l}
    \toprule
    文件 & 内容 \\
    \midrule
    \code{.amadeus/preferences.json} & 偏好库：偏好条目及其版本链、场景标签 \\
    \code{.amadeus/structured\_memory.json} & 知识库：事实版本链、片段、关联边、模板 \\
    \code{.amadeus/mid\_term\_memory.json} & 中期记忆库：网关产出的五类记录 \\
    \code{.amadeus/memory.json} & 长期用户事实（MemoryRegistry 数据源） \\
    \code{.amadeus/rag\_index.json} & 语义向量索引（v2 信封格式，支持 int8 量化） \\
    \bottomrule
  \end{tabular}
\end{table}

\section{核心算法设计}\label{sec:algo}

\subsection{多源偏好记忆：提取、版本化与跨场景适配}\label{sec:pref}

\subsubsection{数据模型}

偏好条目 \code{StoredPreference} 由键、类别、置信度、场景标签与\textbf{版本链}构成。
类别覆盖赛题要求的三类偏好并扩展为六类（表~\ref{tab:prefcat}）。

\begin{table}[htbp]
  \centering\small
  \caption{偏好类别划分}
  \label{tab:prefcat}
  \begin{tabular}{l l}
    \toprule
    类别 & 语义 \\
    \midrule
    \code{ToolChoice} & 工具选择偏好（同任务下倾向使用哪个工具） \\
    \code{OutputStyle} & 输出风格偏好（格式、详略、语言习惯） \\
    \code{SafetyPolicy} & 安全策略偏好（敏感操作授权习惯、避免清单） \\
    \code{Workflow} & 工作流偏好（常用步骤顺序） \\
    \code{Habit} & 重复操作习惯 \\
    \code{General} & 其他通用偏好 \\
    \bottomrule
  \end{tabular}
\end{table}

\subsubsection{三路提取器}

针对赛题要求的三类数据源分别实现确定性提取器（表~\ref{tab:extractors}），全部为
规则与统计阈值驱动，无 LLM 调用。

\begin{table}[htbp]
  \centering\small
  \caption{多源偏好提取器}
  \label{tab:extractors}
  \begin{tabular}{l p{6.6cm} l}
    \toprule
    数据源 & 触发条件与产出 & 产出类别 \\
    \midrule
    工具执行结果 &
      按（任务，工具）聚簇统计成功 / 失败 / 延迟；样本数 $\ge 2$、成功率
      $\ge 0.8$ 且使用占比 $\ge 0.6$ 时生成 \code{tool\_choice::\{task\}} 偏好；
      同簇连续失败 $\ge 2$ 次则对该工具降权（置信度 $-0.15$，不删除条目） &
      ToolChoice \\
    用户行为数据 &
      同一动作重复 $\ge 2$ 次沉淀为 \code{habit::} 偏好；检测到用户纠正行为时生成
      \code{avoid::} 偏好并削弱与之冲突的习惯；安全敏感动作归入安全策略 &
      Habit / SafetyPolicy \\
    手动配置信息 &
      配置项直接生成 \code{config::} 偏好（置信度 $0.98$），按键名关键词推断类别 &
      全类别 \\
    \bottomrule
  \end{tabular}
\end{table}

置信度由样本计数与信号一致性映射：样本越充分、信号越一致，置信度越高；连续失败与
用户纠正均以固定步长衰减，使偏好可随行为漂移而退化。

\subsubsection{版本化管理}

版本链采用\textbf{只追加}策略：任何更新都生成新版本而非改写历史。\code{push\_version}
负责追加版本，历史版本完整保留；\code{rollback} 是非破坏性回滚——生成一个「恢复为
旧值」的新版本而非删除；\code{versions()} 提供完整版本史查询；\code{resolve\_as\_of(t)}
则按任意时间点回溯当时生效的版本，支持审计与复盘。

\subsubsection{跨场景适配与回溯}

每个偏好可携带场景标签（\code{ScenarioTag}：场景名 + 属性键值对）。解析当前生效
偏好时按两级匹配：第一级由 \code{resolve\_for\_scenario} 优先取场景完全匹配的版本，
无则回退默认场景；第二级由 \code{cross\_scenario\_candidates} 按场景属性集合的
Jaccard 相似度排序，跨场景推荐最相近场景下已验证的偏好，实现「跨场景复用」。

\subsubsection{接入状态}

工具执行结果管道已接入 agent 主循环（每次工具调用后即时提取）；用户行为与手动配置
两路提取器已实现，生产接线 \todo{用户行为 / 手动配置管道接入 agent 主循环}；学到的
偏好注入提示词 \todo{偏好注入系统提示词（闭环演示）}。

\subsection{知识记忆：结构化整合、冲突处理与关联检索}\label{sec:knowledge}

\subsubsection{数据模型}

知识库（\code{memory-service}）由四类要素构成。\code{FactVersion} 是 SPO 三元组
（主语—谓语—宾语）附加上下文、版本号与 \code{supersedes\_id} 指针，构成可追溯的
版本链；\code{Episode} 为历史片段，含内容哈希与敏感级别，是模板证据的唯一合法来源；
\code{AssociationEdge} 表示知识条目间的关联边，参与检索加分；\code{MemoryTemplate}
承载可复用工作流模板（前置条件、步骤、输入 schema、证据链、成功率、版本、状态）。

\subsubsection{新旧知识冲突处理}\label{sec:conflict}

新事实入库时先做\textbf{中文同义归一化}（\code{canonicalize}，如「住址 / 地址」统一
为 \code{address}），再与既有事实判定冲突关系，按表~\ref{tab:conflict} 六种动作处置。

\begin{table}[htbp]
  \centering\small
  \caption{冲突判定与六种处置动作}
  \label{tab:conflict}
  \begin{tabular}{l p{3.6cm} p{7.2cm}}
    \toprule
    动作 & 触发条件 & 处置 \\
    \midrule
    \code{Add} & 无既有相关事实 & 新增版本链 \\
    \code{Reinforce} & 同 SPO 且上下文一致 & 置信度 $+0.05$（设上限），不新增版本 \\
    \code{Merge} & 宾语互含（语义包含） & 保留信息量更大（更长）的宾语 \\
    \code{Supersede} & 互斥谓词（住址、工作单位、电话、邮箱、状态等） &
      旧版本转 \code{Historical} 并记录失效时间，新版本接替 \\
    \code{KeepBoth} & 有冲突但无法判定 & 并存，等待后续证据 \\
    \code{Reject} & 输入校验失败（含敏感值、schema 非法） & 拒绝入库 \\
    \bottomrule
  \end{tabular}
\end{table}

被取代的旧事实不删除，转为 \code{Historical} 状态并记录 \code{valid\_to} 时间；
\code{fact\_history()} 可回放任一事实的完整演变史，与偏好版本链的设计一致。

\subsubsection{关联检索}

检索采用「字段匹配 + 关键词 + 关联图 + 状态」的加权打分：

\begin{equation}
  \mathrm{score}(q, f) = 4.0\,[\text{字段全匹配}] + 1.0\,n_{\mathrm{kw}}(q,f)
    + 1.5\,n_{\mathrm{assoc}}(f) + 0.5\,[f \in \mathrm{Active}] + w_c \cdot \mathrm{conf}(f)
\end{equation}

其中 $n_{\mathrm{kw}}$ 为查询词命中数，$n_{\mathrm{assoc}}$ 为命中条目经关联边可达的
证据数，$\mathrm{conf}$ 为置信度。查询词同样经过同义归一化，缓解中文同义表述导致的
漏检。语义级检索（向量近邻）由第~\ref{sec:embedding}~节的 RAG 子系统承担，两者互补。

\subsubsection{可复用模板生命周期}

模板沉淀遵循「证据先行、人审把关」原则（表~\ref{tab:template}）：先以
\code{register\_template} 注册为 \code{Draft} 状态；随后通过
\code{record\_template\_result} 挂接真实执行片段作为证据；满足\textbf{证据数 $\ge 5$、
成功率 $\ge 0.8$、人工复核}三条件方可激活为 \code{Active}；证据所依赖的事实被遗忘时
自动级联降级为 \code{Retired}。

\begin{table}[htbp]
  \centering\small
  \caption{模板状态机}
  \label{tab:template}
  \begin{tabular}{l l}
    \toprule
    状态 & 进入条件 \\
    \midrule
    \code{Draft} & \code{register\_template} 注册 \\
    \code{Active} & 证据 $\ge 5$ 且成功率 $\ge 0.8$ 且人工复核通过（需权限审批） \\
    \code{Retired} & 证据级联失效或人工退役 \\
    \bottomrule
  \end{tabular}
\end{table}

\subsection{敏感信息识别与精准遗忘}\label{sec:forget}

\subsubsection{敏感信息识别与过滤}

\code{privacy} crate 实现六类检测器（表~\ref{tab:privacy}），按三级风险分级，命中
内容替换为 \code{[REDACTED:KIND]} 占位符后入库。脱敏发生在\textbf{所有入库边界}：
偏好入库、中期记忆入库、知识库写入，以及服务启动时对磁盘存量数据的全量复扫。

\begin{table}[htbp]
  \centering\small
  \caption{敏感信息检测器}
  \label{tab:privacy}
  \begin{tabular}{l l}
    \toprule
    类型 & 校验手段 \\
    \midrule
    电子邮箱 & 格式匹配 \\
    手机号 & 号段与位数规则 \\
    身份证号 & 18 位 + 校验码验证 \\
    银行卡号 & Luhn 校验 \\
    API 密钥 & 前缀与熵特征 \\
    口令字段 & 字段名启发（password / token 等） \\
    \bottomrule
  \end{tabular}
\end{table}

\subsubsection{两阶段精准遗忘协议}

遗忘采用「解析—预览—确认—复验」两阶段协议（图~\ref{fig:forget}），全程不可逆操作
均需人工确认。

\begin{figure}[htbp]
  \centering
  \begin{tikzpicture}[node distance=6mm and 8mm]
    \node[box] (parse) {意图解析\\ {\scriptsize parse\_forget\_intent}};
    \node[box, right=of parse] (propose) {生成删除预案\\ {\scriptsize propose\_forget：DeletionPlan}};
    \node[proc, right=of propose] (approve) {人工审批\\ {\scriptsize 权限护栏强制确认}};
    \node[box, below=8mm of approve] (confirm) {执行删除\\ {\scriptsize confirm\_forget：哈希校验 + 物理删除}};
    \node[proc, left=of confirm] (verify) {残留复验\\ {\scriptsize 重读磁盘，断言 0 残留}};
    \node[box, left=of verify] (audit) {审计记录\\ {\scriptsize 无内容 AuditRecord}};

    \draw[arrow] (parse) -- (propose);
    \draw[arrow] (propose) -- (approve);
    \draw[arrow] (approve) -- (confirm);
    \draw[arrow] (confirm) -- (verify);
    \draw[arrow] (verify) -- (audit);
  \end{tikzpicture}
  \caption{两阶段精准遗忘协议}
  \label{fig:forget}
\end{figure}

意图解析支持中英文自然语言指令，含两类修饰语法。\textbf{排除}语法如「忘掉我的所有
住址，\textbf{但保留}城市偏好」，排除项不进入删除预案；\textbf{来源限定}语法如
「只忘掉\textbf{来源：}手动配置的」，按数据来源缩小删除范围。

删除预案（\code{DeletionPlan}）含脱敏预览、预案哈希与 5 分钟生存期；确认时校验哈希
防篡改，物理删除事实、关联边与孤儿片段，级联退役受影响模板，随后\textbf{重读磁盘
验证残留数为零}，最后写入不含任何内容线索的审计记录。无法唯一确定删除目标的模糊指令
（如「忘掉一些东西」）会被安全拒绝。\code{forget\_confirm} 与 \code{activate\_template}
操作在权限层强制一次性人工审批。

\subsection{端侧语义检索：可插拔嵌入与轻量向量库}\label{sec:embedding}

\subsubsection{可插拔嵌入后端}

嵌入能力抽象为 \code{Embedder} trait，按配置在三种后端间切换（表~\ref{tab:backends}），
消费方代码不感知差异。

\begin{table}[htbp]
  \centering\small
  \caption{嵌入后端}
  \label{tab:backends}
  \begin{tabular}{l l l}
    \toprule
    后端 & 形态 & 用途 \\
    \midrule
    \code{remote} & OpenAI 兼容 HTTP 服务 & 开发期 / 有网络环境 \\
    \code{local} & 内置哈希嵌入器，零依赖离线运行 & 端侧基线，无网络可用 \\
    \code{kylin} & 银河麒麟 embedding SDK（本机服务形态） & 麒麟环境部署 \\
    \bottomrule
  \end{tabular}
\end{table}

\code{local} 基线采用特征哈希嵌入：拉丁词 + 中文二元组（bigram）分词后，经 FNV-1a
64 位哈希投射到固定维度（默认 384 维），哈希导出符号位，向量 L2 归一化。该方案完全
确定性、离线、无模型文件，适合端侧兜底；其局限是仅捕获词法重叠，语义近邻能力有限
（优化方向见 \S\ref{sec:conclusion}）。

\subsubsection{轻量化向量存储}

向量库采用单文件 JSON 持久化与内存态线性扫描，轻量化体现在三个方面。\textbf{int8
标量量化}将每向量按绝对值最大值缩放为 int8 存储，查询以整数点积累加（i32 精度），
内存占用降为浮点的四分之一，重启后自相似度 $> 0.99$（有测试断言）；\textbf{v2 信封
格式}在文件头记录版本、量化模式与维度，v1 旧格式透明迁移；\textbf{维度防护}则将与
库内维度不一致的写入与查询显式拒绝，避免静默污染。

检索为 $O(n \cdot d)$ 余弦线性扫描 + 全排序取 top-$k$。在端侧典型知识库规模
（$10^3 \sim 10^4$ 条、$d = 384$）下，单次扫描为千万级浮点乘加，理论耗时远低于
赛题 500ms 上限；分规模实测数据（P50 / P95）\todo{检索延迟实测：库规模 1k / 10k / 100k}。

\section{记忆流转机制：与短期、中期记忆的交互逻辑}\label{sec:flow}

本章对应赛题要求（6）：兼容与记忆模块中短期、中期记忆间的数据流转。

\subsection{短期 $\rightarrow$ 中期：压缩退役网关}

会话上下文接近窗口上限时触发压缩（compaction）。被退役的消息并不直接丢弃，而是经过
\textbf{上下文网关}（\code{ContextGate}）抽取为中期记忆记录（图~\ref{fig:tierflow}）。
网关的确定性实现 \code{RuleBasedGate} 按表~\ref{tab:gate} 的规则识别五类记录。

\begin{table}[htbp]
  \centering\small
  \caption{ContextGate 抽取规则}
  \label{tab:gate}
  \begin{tabular}{l l l}
    \toprule
    退役消息来源 & 记录类型 & 触发条件 \\
    \midrule
    用户文本消息 & \code{Task} & 有效内容 $\ge 8$ 字符，首句作为任务目标 \\
    助手文本消息 & \code{Decision} & 内容含决策关键词（决定 / decision 等） \\
    工具调用入参 & \code{FileState} & 入参含路径型字段（path 等） \\
    工具执行结果 & \code{ErrorResolution} & 结果内容含错误关键词 \\
    压缩摘要 & \code{StateSnapshot} & 压缩发生即产出 \\
    \bottomrule
  \end{tabular}
\end{table}

记录重要性按类型赋默认值（表~\ref{tab:importance}），与压缩摘要策略的优先级一致；
同一次转换内按键去重（后写覆盖），单次产出上限 64 条，防止压缩风暴撑爆存储。

\begin{table}[htbp]
  \centering\small
  \caption{记录类型默认重要性（0--100）}
  \label{tab:importance}
  \begin{tabular}{l c}
    \toprule
    记录类型 & 默认重要性 \\
    \midrule
    \code{Decision} & 80 \\
    \code{StateSnapshot} & 75 \\
    \code{ErrorResolution} & 70 \\
    \code{Task} & 60 \\
    \code{FileState} / \code{Fact} & 50 / 40 \\
    \bottomrule
  \end{tabular}
\end{table}

\subsection{中期记忆存储接口}

中期记忆库实现 \code{MidTermStore} 接口：\code{upsert}（按键替换）、\code{get}、
\code{delete}、\code{query}（多条件 AND 过滤，按时间新到旧排序）、\code{count}、
\code{flush}。记录模型 \code{MemoryRecord} 携带会话来源、消息区间、访问计数与重要性，
可无损转换为长期记忆的 \code{MemoryEntry} 供既有消费方读取。

\subsection{长期记忆注入短期上下文}

长期记忆经 \code{MemoryRegistry}（用户事实提供者 + 向量记忆提供者）在每轮对话构建
系统提示词时注入「Persistent Memory」区块，实现长短期记忆闭环（图~\ref{fig:tierflow}
右侧回边）。

\subsection{中期 $\rightarrow$ 长期晋升}

设计上，中期记录应按「重要性 + 访问计数」驱动的晋升机制进入长期记忆库（重要性衰减
与淘汰同理）。当前实现状态如下：短期 $\rightarrow$ 中期、长期 $\rightarrow$ 短期两条
路径 \statusdone；中期 $\rightarrow$ 长期晋升与中期淘汰 \todo{晋升 / 淘汰机制实现}；
中期记忆查询接入提示词注入 \todo{中期记忆读路径接入 MemoryRegistry}。

\subsection{偏好与知识管道的流转关系}

偏好库与知识库不经过中期库，而是由 agent 主循环直接驱动：工具每次执行后，结果即时
流入偏好提取器；agent 亦可随时经 \code{kylin\_memory} 工具向知识库写入事实、检索、
沉淀模板。两条管道均在入库边界完成脱敏（\S\ref{sec:forget}），与三层流转共用同一
隐私基础设施。

\section{技术实现方案}\label{sec:impl}

\subsection{实现栈与工程约束}

系统以 Rust 实现，全异步 Tokio 运行时，序列化统一为 serde JSON。工程约束如下：全链路
统一错误类型 \code{Result<T>}，生产代码零 \code{unwrap}；公共接口全部 rustdoc 文档化，
模块有文档级说明，满足赛题源代码可读性要求；工作区按 crate 划分（表~\ref{tab:crates}），
记忆子系统可整体独立复用于其他 Agent 产品；特性开关（feature flags）按需裁剪，端侧可
只编译记忆内核 + local 嵌入后端。

\subsection{Agent 主循环接线点}\label{sec:wiring}

记忆系统经统一接线模块 \code{memory\_wiring} 接入 agent 主循环（表~\ref{tab:wiring}），
对主循环零侵入。

\begin{table}[htbp]
  \centering\footnotesize
  \caption{主循环接线点}
  \label{tab:wiring}
  \begin{tabular}{l l p{5.2cm}}
    \toprule
    时机 & 管道 & 目标 \\
    \midrule
    每次工具调用后 & \code{ingest\_tool\_execution} & 偏好库（工具选择偏好提取） \\
    上下文压缩退役时 & \code{archive\_compacted\_context} & 中期记忆库（网关抽取） \\
    系统提示词构建时 & \code{MemoryRegistry} & 长期记忆注入（Persistent Memory 区块） \\
    会话启动时 & 工具注册 & \code{kylin\_memory} / \code{RagTool} 挂载 \\
    服务启动时 & 存量复扫 & 磁盘记忆全量脱敏校验 \\
    \bottomrule
  \end{tabular}
\end{table}

\subsection{知识记忆工具面}

agent 通过 \code{kylin\_memory} 工具操作知识库，九个操作覆盖写入、检索、模板与遗忘
全流程（表~\ref{tab:ops}）；其中 \code{forget\_confirm} 与 \code{activate\_template}
受权限层强制人工审批。

\begin{table}[htbp]
  \centering\small
  \caption{kylin\_memory 工具操作}
  \label{tab:ops}
  \begin{tabular}{l l}
    \toprule
    操作 & 功能 \\
    \midrule
    \code{store\_fact} & 写入 SPO 事实（触发冲突判定与脱敏校验） \\
    \code{search} & 关联检索（加权打分排序） \\
    \code{register\_template} & 注册工作流模板（Draft 状态） \\
    \code{record\_template\_result} & 为模板挂接真实执行证据 \\
    \code{activate\_template} & 激活模板（三条件门槛 + 人工审批） \\
    \code{find\_templates} & 检索可用模板 \\
    \code{privacy\_scan} & 敏感信息扫描 \\
    \code{forget\_preview} & 遗忘预案预览（DeletionPlan） \\
    \code{forget\_confirm} & 确认执行遗忘（哈希校验 + 残留复验） \\
    \bottomrule
  \end{tabular}
\end{table}

\subsection{配置面}

记忆与检索行为均由工作区配置驱动（表~\ref{tab:config}），无需改码即可切换端侧 /
云端 / 麒麟形态。

\begin{table}[htbp]
  \centering\footnotesize
  \caption{关键配置项}
  \label{tab:config}
  \begin{tabular}{l l p{4.4cm}}
    \toprule
    配置项 & 默认值 & 说明 \\
    \midrule
    \code{rag\_enabled} & \code{false} & 是否启用 RAG 工具 \\
    \code{embedding\_backend} & \code{remote} & \code{remote} / \code{local} / \code{kylin} \\
    \code{embedding\_model} & --- & 嵌入模型名 \\
    \code{embedding\_base\_url} & --- & remote 后端服务地址 \\
    \code{kylin\_embedding\_endpoint} & \code{http://127.0.0.1:18080/v1} & 麒麟 SDK 服务端点 \\
    \code{embedding\_dimension} & \code{384} & 向量维度（库内强制一致） \\
    \code{rag\_quantization} & \code{none} & \code{none} / \code{int8} \\
    \code{rag\_chunk\_size} / \code{rag\_chunk\_overlap} & --- & 切块参数 \\
    \code{rag\_top\_k} & --- & 检索返回条数 \\
    \bottomrule
  \end{tabular}
\end{table}

\subsection{质量保障}

记忆子系统建立了三层测试：模块单元测试（知识记忆服务 14 个、偏好提取与版本化 25 个、
隐私检测 3 个等）、跨 crate 集成指标测试（11 个，\S\ref{sec:functest}）、以及全工作区
回归（\code{cargo test --features full}）。2026-09-13 集成指标测试运行结果为
\textbf{11 通过 / 0 失败}。

\section{端侧部署与银河麒麟适配}\label{sec:deploy}

\subsection{端侧轻量化设计}

端侧轻量化体现在四个方面。\textbf{零外部服务依赖}：除可选的嵌入服务外，系统不需要
数据库、消息队列或 GPU，local 嵌入后端零模型文件、完全离线。\textbf{轻量存储}：全部
状态为单文件 JSON + int8 量化向量，内存与磁盘占用小，人可直接检视与备份。
\textbf{确定性算法}：偏好提取 / 冲突处理 / 遗忘判定均为规则与统计阈值，无常驻模型
推理开销。\textbf{按需裁剪}：feature flags 支持只保留记忆内核的最小端侧构建。

\subsection{银河麒麟 embedding SDK 接入}\label{sec:kylin-sdk}

\code{kylin} 嵌入后端当前按「麒麟 SDK 以本机 OpenAI 兼容服务进程形态提供」接入
（默认端点 \code{http://127.0.0.1:18080/v1}）。针对 SDK 交付形态的可能变化，系统
预留了无侵入的升级路径（表~\ref{tab:kylinpath}）：由于嵌入能力收敛在 \code{Embedder}
trait 之后，任何形态变化只影响 \code{KylinEmbedder} 单个文件的内部实现。

\begin{table}[htbp]
  \centering\small
  \caption{麒麟 SDK 交付形态与升级路径}
  \label{tab:kylinpath}
  \begin{tabular}{l l}
    \toprule
    SDK 形态 & 适配方式 \\
    \midrule
    本机 HTTP 服务（OpenAI 兼容） & 现有后端直接可用，仅改端点配置 \\
    C FFI 动态库 & 重写 \code{KylinEmbedder} 内部加载逻辑，接口面不变 \\
    嵌入维度不同 & 维度防护显式拒绝混用，重建向量索引后重新注入 \\
    \bottomrule
  \end{tabular}
\end{table}

与麒麟 SDK 的真机联调记录 \todo{麒麟 embedding SDK 真机联调（SDK 到货后执行）}。

\subsection{安装部署指南}

\textbf{前置依赖}：Rust 稳定版工具链；可选的嵌入服务（remote 形态或麒麟 SDK 服务）。

\textbf{构建与运行}：

\begin{lstlisting}[language=bash]
# 构建（无默认特性，按需启用）
cargo build --features full
# 运行（TUI / HTTP 服务）
cargo run --features full
cargo run --features full -- --server 8080
\end{lstlisting}

\textbf{端侧启用麒麟嵌入 + int8 量化的最小配置}（工作区 \code{.amadeus/} 下）：

\begin{lstlisting}[language={}]
{
  "rag_enabled": true,
  "embedding_backend": "kylin",
  "kylin_embedding_endpoint": "http://127.0.0.1:18080/v1",
  "embedding_dimension": 384,
  "rag_quantization": "int8",
  "rag_top_k": 5
}
\end{lstlisting}

运行时目录 \code{.amadeus/}（表~\ref{tab:storage}）随首次运行自动创建。

\subsection{操作说明}

记忆操作经 agent 的 \code{kylin\_memory} 工具以自然语言或结构化指令完成，典型示例：

\begin{lstlisting}[language={}]
{"operation": "store_fact", "subject": "user",
 "predicate": "住址", "object": "南京市", "context": "工作信息"}

{"operation": "search", "query": "住址"}

{"operation": "forget_preview",
 "instruction": "忘掉我的所有住址，但保留城市偏好"}
\end{lstlisting}

遗忘操作先预览预案、再经人工审批确认（\S\ref{sec:forget}）；模板激活同理。

\subsection{银河麒麟桌面操作系统适配测试报告}

记忆子系统全部依赖（Tokio、serde、reqwest 等）均原生支持 Linux / 国产 CPU 架构，
代码中不含平台专有调用，具备跨平台适配基础。正式适配测试记录如表~\ref{tab:kylinenv}。

\begin{table}[htbp]
  \centering\small
  \caption{银河麒麟适配测试环境与结果}
  \label{tab:kylinenv}
  \begin{tabular}{l l}
    \toprule
    项目 & 记录 \\
    \midrule
    操作系统版本 & \todoblank \\
    内核版本 / CPU 架构 & \todoblank \\
    构建工具链版本 & \todoblank \\
    功能测试结果 & \todoblank \\
    性能指标结果 & \todoblank \\
    \bottomrule
  \end{tabular}
\end{table}

\todo{银河麒麟桌面 OS 环境完成安装部署与全量回归，出具适配测试报告}

\section{效果验证报告}\label{sec:evaluation}

\subsection{评测指标定义与口径}\label{sec:metrics}

四项核心指标的评测口径定义如表~\ref{tab:metricdef}，全部指标在标注数据集上以自动化
脚本计算，可复现。

\begin{table}[htbp]
  \centering\small
  \caption{评测指标定义}
  \label{tab:metricdef}
  \begin{tabular}{l c p{7.2cm}}
    \toprule
    指标 & 目标值 & 口径定义 \\
    \midrule
    偏好提取准确率 & $\ge 85\%$ &
      系统提取的偏好与人工标注偏好（类别 + 内容 + 场景）的一致率 \\
    知识检索召回率 & $\ge 85\%$ &
      标注相关的知识条目出现在 top-$k$（$k=5$）结果中的比例（Recall@5） \\
    知识检索响应时间 & $\le 500\mathrm{ms}$ &
      端到端检索耗时（含打分排序）的 P50 / P95 \\
    知识冲突处理正确率 & $\ge 88\%$ &
      六动作判定结果与人工标注期望动作的一致率 \\
    \bottomrule
  \end{tabular}
\end{table}

\subsection{标准化测试数据集说明}\label{sec:dataset}

评测数据集由三个子集构成。\textbf{偏好标注集}收录工具调用执行序列（任务—工具—
成败—延迟）、用户行为序列与配置项，由双人独立标注「应提取出的偏好」并仲裁分歧，
用于偏好提取准确率。\textbf{知识检索集}为查询—相关知识条目对（含同义改写、多关联边、
历史版本干扰项），用于 Recall@5 与检索延迟。\textbf{冲突判定集}为新旧事实对附期望
处置动作（六动作全覆盖，含互斥谓词与上下文一致性边界情形），用于冲突处理正确率。

数据集规模、来源与标注一致性（双人一致率 / Kappa）\todo{评测数据集构建完成并冻结
版本}。

\subsection{量化指标结果}

\begin{table}[htbp]
  \centering\small
  \caption{核心指标实测结果}
  \label{tab:metricresult}
  \begin{tabular}{l c c c}
    \toprule
    指标 & 目标值 & 实测值 & 结论 \\
    \midrule
    偏好提取准确率 & $\ge 85\%$ & \todoblank & \todoblank \\
    知识检索召回率（Recall@5） & $\ge 85\%$ & \todoblank & \todoblank \\
    检索响应时间 P50 / P95 & $\le 500\mathrm{ms}$ & \todoblank & \todoblank \\
    冲突处理正确率 & $\ge 88\%$ & \todoblank & \todoblank \\
    \bottomrule
  \end{tabular}
\end{table}

\todo{在冻结数据集上运行四项指标评测脚本并填入本表；未达标项在第~\ref{sec:conclusion}~
章给出优化方向}

\subsection{对比实验设计（消融）}

为量化各机制的贡献，设计如下消融对比（表~\ref{tab:ablation}），实验组仅关闭单一
机制，其余条件一致。

\begin{table}[htbp]
  \centering\small
  \caption{消融实验设计与结果}
  \label{tab:ablation}
  \begin{tabular}{l l c}
    \toprule
    实验组 & 检验假设 & 结果 \\
    \midrule
    完整方案 & — & \todoblank \\
    去版本化（仅保留最新版本） & 版本链对回溯 / 审计的贡献 & \todoblank \\
    去冲突融合（新事实一律 Add） & 六动作机制对正确率的贡献 & \todoblank \\
    去关联边（检索不加关联分） & 关联图对召回的贡献 & \todoblank \\
    去同义归一化 & 中文同义对召回的贡献 & \todoblank \\
    int8 vs 全精度 & 量化对检索质量与延迟的影响 & \todoblank \\
    \bottomrule
  \end{tabular}
\end{table}

\subsection{功能验收测试}\label{sec:functest}

跨 crate 集成指标测试共 11 个，2026-09-13 以 \code{cargo test --features full} 运行，
\textbf{全部通过}（表~\ref{tab:functest}），覆盖赛题内容要求 (3)(4)(5) 的核心闭环。

\begin{table}[htbp]
  \centering\footnotesize
  \caption{集成指标测试（11 / 11 通过）}
  \label{tab:functest}
  \begin{tabular}{p{6.4cm} p{5.2cm} c}
    \toprule
    测试 & 验证点 & 结果 \\
    \midrule
    \testname{indicator_3_conflict_search_and_history_form_a_closed_loop} &
      冲突接替—检索—版本史闭环 & 通过 \\
    \testname{indicator_3_template_activation_requires_five_examples_review_and_eighty_percent_success} &
      模板激活三条件门槛 & 通过 \\
    \testname{indicator_3_default_agent_registers_kylin_memory_without_legacy_registry} &
      工具默认注册 & 通过 \\
    \testname{indicator_5_rejects_ambiguous_and_non_deletion_requests} &
      模糊遗忘指令安全拒绝 & 通过 \\
    \testname{indicator_5_sensitive_information_is_redacted_at_scan_and_storage_boundaries} &
      扫描与入库双边界脱敏 & 通过 \\
    \testname{indicator_5_precise_forgetting_preserves_exclusions_and_requires_confirmation} &
      排除项保留 + 强制确认 & 通过 \\
    \testname{indicator_4_local_hash_embedder_is_deterministic_normalized_and_offline} &
      local 后端确定性 / 归一化 / 离线 & 通过 \\
    \testname{indicator_4_int8_quantized_store_survives_restart} &
      int8 量化重启自相似度 $> 0.99$ & 通过 \\
    \testname{indicator_4_rag_tool_ingest_and_query_end_to_end} &
      注入—检索端到端 & 通过 \\
    \testname{indicator_4_backend_selection_from_config} & 后端配置切换 & 通过 \\
    \testname{indicator_4_agent_registers_rag_tool_when_enabled} &
      RAG 工具按配置挂载 & 通过 \\
    \bottomrule
  \end{tabular}
\end{table}

模块级单元测试另有数十个（知识记忆服务 14、偏好提取与版本化 25、隐私检测 3 等），
覆盖阈值边界、版本回滚、跨场景解析与校验码算法。

\subsection{历史基线与重测计划}

2026 年 5 月曾以 LoCoMo 长对话数据集（10 会话 / 189 问答对）对早期 RAG 链路做过一次
评测（当时为 remote 嵌入后端、LLM 问答 F1 口径）：RAG 模式 F1 = 39.2\%、全上下文
F1 = 54.0\%，其中单跳 60.4\%、多跳 25.0\%、时序 12.2\%，平均查询耗时 0.02s。

需要说明：\textbf{该口径度量的是问答生成质量而非检索召回率}，且早于 local / kylin
后端与 int8 量化的引入，不能直接对标赛题指标，仅作为演进基线保留。以本方案口径
（Recall@5，local / kylin 后端）的重测 \todo{LoCoMo 数据集按 Recall 口径重测}。

\subsection{检索延迟实测方案}

延迟实测按库规模 $\times$ 量化模式网格执行（1k / 10k / 100k 条，int8 / 全精度，
$d=384$），每组 1000 次查询取 P50 / P95，测试机规格与结果
\todo{延迟网格实测并填入 assets/metrics-latency.png}。

\section{真实场景应用案例}\label{sec:case}

\subsection{场景设计：开发助手日常使用（办公自动化）}

选取「基于 OS Agent 的开发 / 办公自动化助手」作为真实落地场景，覆盖三类记忆能力的
完整闭环。偏好沉淀与复用方面，用户在多轮会话中反复于同类任务选择某一工具、要求简洁
输出格式，系统从工具执行结果与用户行为中自动提取工具选择与输出风格偏好，跨会话、
跨场景（终端场景 / 文档场景）持续生效；用户手动修改配置后，偏好经版本化更新并可
回溯。知识沉淀与智能调用方面，一次完整的「构建—排错—修复」流程被沉淀为可复用模板
（含前置条件与步骤），五次真实执行证据且成功率达标并经人工复核后激活，后续同类故障
一键复用；新旧事实（如项目构建参数变更）经冲突处理自动接替，版本史可查。安全合规
方面，记忆过程中出现的手机号、邮箱等被自动脱敏；演示「忘掉我的所有住址，但保留城市
偏好」——排除项保留、目标数据物理删除、磁盘零残留。

\subsection{实测记录}

案例执行记录、关键步骤截图与启用 / 关闭记忆的前后对比 \todo{真实场景案例执行并
记录（步骤截图、前后对比、用户反馈）}。

效果演示视频 \todo{录制偏好捕捉、知识整合、检索复用三段演示视频}。

\subsection{应用价值}

该场景直接对应赛题「办公自动化、智能协作、个性化助手」的应用方向：偏好记忆减少
重复指令与工具选择成本，知识模板把一次性排错经验转化为可复用资产，精准遗忘满足
个人信息保护的合规要求。量化收益（任务完成时间缩短比、指令次数下降比）
\todo{案例量化收益统计}。

\section{总结与优化方向}\label{sec:conclusion}

\subsection{已完成工作}

截至本报告撰写时，系统已完成三层记忆架构（短期 / 中期 / 长期）与双库（偏好 / 知识）
落地，全部状态文件化、零外部服务依赖；实现多源偏好提取三路提取器、只追加版本化、
非破坏回滚、按时间点回溯与基于场景相似度的跨场景复用；建成 SPO 事实版本链、同义
归一化、六动作冲突处理、关联加权检索与证据先行的模板生命周期；落实六类敏感信息检测、
全入库边界脱敏与两阶段精准遗忘（预案—审批—物理删除—残留复验—审计）；交付可插拔
嵌入三后端（remote / local / kylin）、int8 量化向量库与 v1/v2 格式迁移；记忆子系统
已接入 agent 主循环，11 个集成指标测试与数十个模块单元测试通过（2026-09-13 运行）。

\subsection{待完成事项}

以下事项均在前文以红色 TODO 标注，此处按工作类别汇总。（1）\todo{标准化评测数据集
三子集构建与冻结（\S\ref{sec:dataset}）}；（2）\todo{四项性能指标实测与消融实验
（\S\ref{sec:metrics}）}；（3）\todo{检索延迟网格实测（\S\ref{sec:metrics}）}；
（4）\todo{用户行为 / 手动配置管道接入 agent 主循环（\S\ref{sec:pref}）}；
（5）\todo{偏好注入系统提示词，形成个性化闭环演示（\S\ref{sec:pref}）}；
（6）\todo{中期 $\rightarrow$ 长期晋升与中期淘汰机制（\S\ref{sec:flow}）}；
（7）\todo{麒麟 embedding SDK 真机联调（\S\ref{sec:kylin-sdk}）}；（8）\todo{银河
麒麟桌面 OS 适配测试报告（\S\ref{sec:deploy}）}；（9）\todo{真实场景案例实测记录与
演示视频（\S\ref{sec:case}）}；（10）\todo{LoCoMo 数据集按 Recall 口径重测
（\S\ref{sec:evaluation}）}。

\subsection{后续优化方向}

后续优化沿四个方向推进。\textbf{本地神经嵌入}：以小型中文嵌入模型（int8 量化）替换
哈希基线，提升语义召回上限，同时保持离线与低资源占用。\textbf{混合检索排序}：结构化
加权检索与向量近邻的融合排序，兼顾精确字段匹配与语义泛化，知识库规模增长后引入轻量
ANN 索引。\textbf{遗忘意图的 NLU 扩展}：在保留规则式安全护栏的前提下，扩展意图解析
的概念覆盖面。\textbf{中期记忆生命周期}：重要性衰减、访问计数驱动晋升与淘汰，补全
三层流转的最后一环。


\end{document}
